Com o crescimento de aplicações baseadas em grandes volumes de dados, a
necessidade de processar relações entre entidades tornou-se cada vez mais
relevante. Uma forma natural de representar essas relações é por meio de grafos,
nos quais vértices representam entidades e arestas representam conexões entre
elas.

Em uma rede social, vértices podem representar usuários e arestas podem
representar amizades, seguidores, mensagens ou interações \cite{easleykleinberg}.
Na Web, páginas são vértices e hiperlinks são arestas \cite{brinpage}. Em
segurança computacional, grafos podem modelar fluxos de rede, dependências entre
processos, tentativas de acesso e relações entre entidades suspeitas
\cite{akoglugraphanomaly}. Em biologia, vértices podem representar proteínas,
genes ou moléculas, e arestas podem representar interações
\cite{barabasinetworkbiology}. Em todos esses casos, a Busca em Largura, ou BFS
(\textit{Breadth-First Search}), aparece como uma primitiva para calcular
alcançabilidade, distância em número de arestas e expansão de vizinhança.

Este trabalho investiga a execução de BFS em GPU usando CUDA e NVSHMEM. O
Graph500 foi utilizado como benchmark de referência, pois sua especificação
avalia justamente cargas de trabalho centradas em grafos, incluindo geração de
grafos sintéticos, construção de estruturas internas e execução de busca em
largura. A partir desse benchmark, foram desenvolvidas duas etapas principais.

Primeiro, foi implementada uma versão CUDA/NVSHMEM dos kernels do Graph500, com
foco no Kernel 1, responsável pela construção da representação do grafo, e no
Kernel 2, responsável pela BFS. Essa versão mantém compatibilidade conceitual com
o fluxo do Graph500, mas executa a expansão das fronteiras em GPU e utiliza
NVSHMEM para comunicação entre elementos de processamento.

Segundo, a implementação foi refatorada em direção a uma biblioteca de travessia
de grafos executada majoritariamente em GPU. Nessa versão, a geração do grafo, a
conversão para uma representação distribuída e a execução da BFS são organizadas
como componentes separados. A principal mudança em relação à versão híbrida é que
o laço principal da BFS passa a permanecer no dispositivo, reduzindo a
participação recorrente da CPU entre níveis da busca.

Os experimentos foram realizados em dois ambientes. O primeiro foi o computador
pessoal usado no desenvolvimento, com processador Ryzen 5 5600X, 32 GB de memória
RAM e uma GPU NVIDIA RTX 5060 Ti com 16 GB de memória. O segundo foi um ambiente
GB200 com processador Grace, 128 threads de CPU e quatro GPUs GB200.

Os resultados mostram que a geração R-MAT em GPU apresenta ganhos expressivos nas
maiores escalas avaliadas. No computador pessoal, em \texttt{SCALE=22}, a geração
em GPU foi aproximadamente $7{,}4$ vezes mais rápida que a geração da referência
em CPU. Para o Kernel 2, a versão CUDA/NVSHMEM híbrida superou a referência em
CPU a partir de \texttt{SCALE=21}. A biblioteca GPU apresentou ganhos ainda mais
expressivos: no computador pessoal, alcançou 2,138 GTEPS em \texttt{SCALE=21},
enquanto a versão híbrida alcançou 0,236604 GTEPS na mesma escala. No ambiente
GB200, em \texttt{SCALE=29}, a biblioteca GPU alcançou 6,180 GTEPS, contra 1,570
GTEPS da referência em CPU, além de reduzir o tempo de geração do grafo de 31,13
s para 2,738 s.

Assim, este trabalho contribui com uma análise da BFS em GPU usando CUDA e
NVSHMEM, uma implementação compatível com o fluxo do Graph500 e uma versão
refatorada em biblioteca que reduz a dependência do host durante a travessia. Os
resultados indicam que mover a coordenação da BFS para a GPU é uma direção
promissora para travessia distribuída de grafos, embora o desempenho continue
dependendo de fatores como irregularidade do grafo, comunicação remota,
sincronização e uso eficiente da memória.
